\section{Implementation}
\label{sec:implementation}

The primitive lives in three Rust components: a runtime kernel that exposes a pre-dispatch admission hook, a federation crate that implements the strict bilateral DSSE verifier, and a selective-disclosure module that projects canonical receipt bytes into a BBS message vector for presentation. The wire format is the load-bearing artifact; the implementation is one realization of it.

\paragraph{Pre-dispatch admission hook.} The runtime exposes a single hook that intercepts every tool-call request before dispatch. Non-federated local requests use the existing local-policy path when no Chiodos admission context is present; once a Chiodos context appears, malformed context denies, and federated origin without treaty context denies with a named rejection code. The hook resolves the treaty reference from the call site rather than from request metadata: request-level smuggling of trust roots, peer directories, signing keys, or dynamic trust bundles is rejected by construction, because the resolver consults only verifier-owned state. Treaty DSSE evidence is verified against verifier-owned trust store artifacts; signer sets must equal the treaty's declared participants and public keys must be independent. The hook releases reserved continuation state on denial or abort, closing the replay edge that an attacker would otherwise exploit by re-submitting a continuation that was admitted-then-aborted.

The asymmetry is deliberate. Allowing federated origin without treaty context would make the bilateral DSSE primitive irrelevant to cross-kernel execution. Allowing request metadata to smuggle a trust root would collapse the verifier-owned-store assumption that the freestanding theorem of \S\ref{sec:formal} rests on. Failing to release reserved continuation state on denial would create a state leak and an ambiguous replay surface. Each is a concrete place where agent middleware turns formal policy into advisory logging; the hook refuses each pattern.

\paragraph{Strict bilateral verifier.} The federation crate at \codepath{crates/chio-federation/src/bilateral_dsse.rs} defines the strict predicate and envelope. The verifier function \codepath{verify_chiodos_dsse_envelope} implements the five-gate check of \S\ref{sec:predicate}: it rejects wrong payload type, wrong signature count, noncanonical JSON, wrong statement type, wrong predicate type, missing subject, and reused signer keys. A companion operational verifier at \codepath{crates/chio-federation/src/bilateral_verifier.rs} layers checks the freestanding theorem treats as trust-store assumptions: peer key pins, fresh ladder manifests, revocation epoch, receipt resolution and signature, tool-argument hash, policy-verdict agreement, capability lease, governance receipt for receipt-backed action classes, and the treaty binding references. Together, the envelope check and the operational check ensure admission depends on the predicate the protocol names, not on a strict subset a sender could satisfy with degenerate data.

A DSSE envelope with the wrong predicate is not a near miss; a payload signed by the same key twice is not two-party consent; a receipt with a valid signature but the wrong request hash is not admissible under the treaty; receipt-backed classes require a governance receipt even when the ordinary receipt is otherwise valid. The five-code error taxonomy is part of the protocol because downstream audit must know which constitutional precondition failed.

\paragraph{Rejection codes mapped to source.} Table~\ref{tab:rejection-codes} maps each of the five rejection codes of \S\ref{sec:predicate} to the verifier source location of the gate that returns the code. The envelope check and the operational check are layered: the envelope check gates on structural and signer-independence properties of the DSSE bytes, while the operational check gates on verifier-owned state (revocation epoch, capability lease, treaty bindings) only on inputs that have already passed the envelope check.

\begin{table}[t]
\footnotesize
\caption{Five rejection codes mapped to verifier source locations.}
\label{tab:rejection-codes}
\begin{tabularx}{\columnwidth}{@{}lX@{}}
\toprule
Rejection code & Verifier source \\
\midrule
noncanonical-payload & canonical-bytes equality gate inside \codepath{verify_chiodos_dsse_envelope} (\codepath{bilateral_dsse.rs:996}) \\
predicate-type-mismatch & \codepath{predicate_type} equality against \codepath{PREDICATE_TYPE_CHIODOS_BILATERAL} inside \codepath{verify_chiodos_dsse_envelope} (\codepath{bilateral_dsse.rs:996}) \\
signer-reuse & keyid-independence gate and \codepath{require_unique_signature_keyids} inside \codepath{verify_chiodos_dsse_envelope} (\codepath{bilateral_dsse.rs:996}) \\
stale-lease & capability-lease freshness gate inside the operational verifier (\codepath{bilateral_verifier.rs}) \\
subject-digest-mismatch & binding-tuple subject equality gate inside the operational verifier (\codepath{bilateral_verifier.rs}) \\
\bottomrule
\end{tabularx}
\end{table}

\paragraph{Selective disclosure.} The selective-disclosure module signs a BBS projection of the canonical receipt body and verifies derived disclosure proofs against an issuer-owned key registry~\cite{boneh2004bbs,bbsDraft}. The Ed25519 signature over canonical receipt bytes remains authoritative for audit corpora; BBS is a presentation-only commitment that lets a holder reveal a subset of receipt fields to a verifier without exposing every field. The verifier pins the issuer key and the projection version, so the proof is not free-form redaction.

\paragraph{Stub disclosure.} One platform-specific stub the primitive does not depend on is worth naming. The macOS Endpoint Security path used for kernel-side telemetry is a stub: the integration surface compiles and tests pass against canned fixtures, but live Endpoint Security event delivery is not wired up. The bilateral-DSSE primitive, the admission hook, the federation crate, and the selective-disclosure module operate on canonical bytes and trust-store state regardless of whether kernel-side telemetry is live; the stub is named so a deployer planning kernel-event-driven receipts knows where the gap sits.
